\documentclass[12pt,a4paper]{article}
\usepackage[margin=1in]{geometry}
\usepackage{amsmath,amssymb,amsthm}
\usepackage{hyperref}
\usepackage{enumitem}

\newtheorem{theorem}{Theorem}[section]
\newtheorem{corollary}[theorem]{Corollary}
\newtheorem{lemma}[theorem]{Lemma}
\newtheorem{definition}[theorem]{Definition}
\newtheorem*{remark}{Remark}

\title{\textbf{The Cascade Series --- Part IVb}\\[6pt]
The Standard Model from the Cascade:\\
Masses, Couplings, and Precision Predictions\\
from the Geometric-Topological Factorization}
\author{RTAC}
\date{March 2026}

\begin{document}
\maketitle

\begin{abstract}
The cascade series tests one hypothesis: the infinite-dimensional unit ball, descended to four dimensions, is indistinguishable from our universe. The companion paper [4] derives the Standard Model gauge group $\text{SU}(3) \times \text{SU}(2) \times \text{U}(1)$, its symmetry breaking pattern, and three fermion generations from the cascade's geometry. This paper derives the fermion mass spectrum, gauge coupling constants, and nine independent precision predictions with zero free parameters.

The central new result is the geometric-topological factorization of the fermion mass. Every particle is a feature of the cascade's higher-dimensional geometry. Its projection into 4D spacetime is attenuated by two independent channels: (1) a geometric channel $\exp(-\Phi(d_g))$, governed by the digamma function, encoding the distance of the feature from the observer; and (2) a topological channel $(2\sqrt{\pi})^{-n_D}$, governed by $\Gamma(\tfrac{1}{2})$ and the Euler characteristic $\chi(S^{2n}) = 2$, encoding the number of hairy ball obstructions between the feature and the observer. These factorize because topological invariants do not depend on the metric.

The topological obstruction factor is $2\sqrt{\pi} = 2\,\Gamma(\tfrac{1}{2})$ per Dirac layer: 2 from chirality ($\chi(S^{2n}) = 2$) and $\sqrt{\pi} = \Gamma(\tfrac{1}{2})$ from the cascade's universal quarter-turn constant. The $d$-dependent coupling integral $R'(d)$ does not appear---confirmed by decisive numerical exclusion (61--69\% deviation with $R'(d)$ vs 0.13--1.7\% without). The universal coupling $C = \alpha_s/(2\sqrt{\pi})$ eliminates the last free parameter: the observer at $d = 4$ sits behind the $d = 5$ hairy ball zero and pays one topological toll to see the gauge window.

The complete mass formula
\[
m_g = \frac{\alpha_s \cdot v}{\sqrt{2}}\,\exp\!\bigl(-\Phi(d_g)\bigr)\,(2\sqrt{\pi})^{-(n_D+1)}
\]
predicts: $m_\tau = 1755$~MeV (obs 1777, 1.2\%), $m_\mu = 106.2$~MeV (obs 105.7, 0.47\%), $m_e = 0.514$~MeV (obs 0.511, 0.60\%). Combined with the gauge coupling predictions ($1/\alpha_{\text{GUT}} = 25.02$, $\alpha_s = 0.1159$, $\sin^2\theta_W = 0.2325$, $m_H = 125.49$~GeV), and the cosmological matter fraction $\Omega_m = 0.31150$ from the Bott partition of cascade layers (Papers I and IVa), the series produces ten independent and four derived precision predictions, all sub-3\%, with zero free parameters.

The cascade's predictions separate into two populations with distinct error signatures. Descent-dependent quantities ($\alpha_s$, $v$, $\Omega_m$, mass ratios, $m_\tau$, $m_W$) carry uniformly negative deviations of $-0.13\%$ to $-2.20\%$ (7/7 negative). Geometric quantities ($m_H/m_W$, $\sin^2\theta_W$, $\theta_C$, $\Omega_b$) carry positive deviations of $+0.56\%$ to $+2.76\%$ (4/4 positive). The sign separation identifies the subleading cascade descent as the source of the dominant systematic.
\end{abstract}

\tableofcontents
\newpage

\section{Division of Labour}

Part IVa [4] derives the Standard Model's structure: the gauge group $\text{SU}(3) \times \text{SU}(2) \times \text{U}(1)$ from the Bott mirror, the symmetry breaking pattern from the hairy ball theorem, and three generations from the $d_1 = 19$ phase transition. This paper derives the Standard Model's quantities: masses, couplings, and precision predictions.

\section{The Fermion Mass Spectrum}

\subsection{The cascade potential}

\begin{definition}[Cascade potential]
The cascade potential at dimension $d$ is
\[
\Phi(d) := \sum_{d'=5}^{d} p(d') = \sum_{d'=5}^{d} \Bigl[\tfrac{1}{2}\psi\!\Bigl(\frac{d'+1}{2}\Bigr) - \tfrac{1}{2}\ln\pi\Bigr],
\]
where $\psi = \Gamma'/\Gamma$ is the digamma function.
\end{definition}

The cascade potential is the cumulative decay rate, measuring the total suppression of a propagator traversing from $d = 4$ to dimension $d$. At the generation layers: $\Phi(5) = -0.111$, $\Phi(13) = 1.429$, $\Phi(21) = 5.494$, $\Phi(29) = 11.082$.

\subsection{The geometric-topological factorization}

Every particle is a feature of the cascade at some dimension $d$ in the higher-dimensional geometry. The 4D observer sees its projection into 4D, attenuated by every layer it passes through. This projection decomposes into two independent channels.

\textbf{Geometric channel.} The cascade potential $\Phi(d_g)$ measures the cumulative suppression through the cascade's sphere-area decay. The geometric attenuation is $\exp(-\Phi(d_g))$. It is continuous, $d$-dependent, and governed by the digamma function.

\textbf{Topological channel.} Between $d = 4$ and the generation layer $d_g$, the projection passes through $n_D(d_g)$ Dirac layers. At each Dirac layer $d$ (where $d \bmod 8 = 5$), the sphere $S^{d-1}$ is even-dimensional and the hairy ball theorem forces a zero. Each obstruction attenuates the projection by $2\sqrt{\pi}$, giving $(2\sqrt{\pi})^{-n_D}$.

The two channels factorize because topological invariants do not depend on the metric.

\subsection{The topological obstruction factor}

\begin{theorem}[Topological obstruction factor]\label{thm:obstruction}
The fermion-to-scalar propagator ratio at each Dirac layer is $1/(2\sqrt{\pi})$, independent of $d$.
\end{theorem}

\begin{proof}
The cascade lapse at layer $d$ factorises as $N(d) = \sqrt{\pi}\cdot R(d)$
(Part~0, Theorem~3.1), where $\sqrt{\pi} = \Gamma(\tfrac{1}{2})$ is the universal
compression constant from orthogonality and
$R(d) = \Gamma((d+1)/2)/\Gamma((d+2)/2)$ is the $d$-dependent geometric coupling.
At a Dirac layer ($d\bmod 8 = 5$), $S^{d-1}$ is even-dimensional.
Two independent obstructions attenuate the fermion propagator:

\medskip
\textbf{(a) The chirality obstruction (factor $\mathbf{1/2}$).}
On an even-dimensional sphere $S^{2n}$, the spinor bundle decomposes as
$S = S^+\oplus S^-$ under the chirality grading. The Euler characteristic
$\chi(S^{2n}) = 2$ counts the minimum number of critical points of any Morse
function on $S^{2n}$ (Poincar\'{e}--Hopf theorem applied to its gradient).
These two critical points---one source, one sink of the height function's
gradient flow---define two chirality basins. A definite-chirality fermion
occupies one basin, halving the propagator amplitude. The factor
$1/\chi = 1/2$ is dimension-independent because $\chi(S^{2n}) = 2$ for all
$n\geq 1$.

\medskip
\textbf{(b) The quarter-turn obstruction (factor $\mathbf{1/\sqrt{\pi}}$).}
The cascade's universal constant $\sqrt{\pi}$ enters the scalar propagator
through the half-integer first argument of $B(\tfrac{1}{2},\cdot)$. This
half-integer arises because the orthogonal quarter-turn from slicing axis to
equator produces a Jacobian $t^{-1/2}$ (under $t = x^2$) in the Beta integral:
\[
B(\tfrac{1}{2},\, d/2+1) = \int_0^1 t^{-1/2}(1-t)^{d/2}\,dt
= \sqrt{\pi}\cdot R'(d).
\]
The $t^{-1/2}$ factor encodes the angular spread of the tangent frame as the
slicing axis rotates from pole to equator. It is a \emph{frame-dependent}
quantity: it measures how the tangent directions fan out along the quarter-turn.

On an odd-dimensional sphere $S^{2n+1}$, a nowhere-vanishing tangent vector
field exists (the hairy ball theorem does not apply). The tangent frame
completes the quarter-turn smoothly, and $\sqrt{\pi}$ enters the propagator
in full.

On an even-dimensional sphere $S^{2n}$, the hairy ball theorem forces every
continuous tangent vector field to vanish. The tangent frame is globally
obstructed. A scalar (0-form) couples only to the metric, which is
well-defined everywhere regardless of the tangent frame---it sees the full
quarter-turn measure $\sqrt{\pi}$. A fermion (spinor) couples to the tangent
frame through the spin connection. The global obstruction of the tangent frame
means the fermion cannot access the full angular measure of the quarter-turn.
The quarter-turn constant $\sqrt{\pi}$ is consumed by the obstruction.

The geometric coupling $R(d)$, which depends on the radial structure of the
slicing (the ratio of adjacent Gamma functions), measures the change in
cross-sectional size---a metric quantity that does not require a tangent frame.
This factor survives for both scalars and fermions.

\medskip
\textbf{(c) Uniqueness.}
The two factors $\sqrt{\pi}$ and $\chi = 2$ are the \emph{only}
dimension-independent constants available at the hairy ball obstruction:
$\sqrt{\pi} = \Gamma(\tfrac{1}{2})$ is the unique constant of the cascade's
slicing recurrence (Part~0, Theorem~4.1); $\chi(S^{2n}) = 2$ is the unique
topological invariant of even-dimensional spheres that counts the obstruction
strength. Any $d$-dependent factor would break the universality of the mass
formula, contradicting the observed $d$-independence of the Bott factor to
1.6\% (Remark~\ref{rem:universality}).

\medskip
The fermion lapse at a Dirac layer is therefore:
\[
N_f(d) = \frac{R(d)}{\chi(S^{d-1})} = \frac{R(d)}{2}.
\]
The fermion-to-scalar ratio is:
\[
\frac{N_f(d)}{N(d)} = \frac{R(d)/2}{\sqrt{\pi}\cdot R(d)} = \frac{1}{2\sqrt{\pi}},
\]
independent of $d$.
\end{proof}

\noindent The factor decomposes as: (a) $\sqrt{\pi} = \Gamma(\tfrac{1}{2})$: the cascade's unique dimension-independent constant (Paper I, Theorem 3.1), forced by orthogonality---consumed at the hairy ball zero because the tangent frame it encodes is globally obstructed. (b) $2 = \chi(S^{2n})$: the Euler characteristic of any even-dimensional sphere, equal to the number of chirality basins at the Dirac junction.

\subsection{Why $R'(d)$ does not appear}

The Higgs coupling integral on $S^{2n}$ at Dirac layer $d$ is $B(\tfrac{1}{2}, \tfrac{d}{2}) = \sqrt{\pi} \times R'(d)$, where $R'(d) = \Gamma(d/2)/\Gamma((d+1)/2)$ is $d$-dependent. This integral computes the geometric overlap of the Higgs field with the fermion wavefunction on a specific sphere. It is the wrong quantity.

The mass coupling at the hairy ball zero is topological: it depends on the existence and index of the zero, not on the geometry of the sphere it lives on. The $d$-dependent factor $R'(d)$ is geometric information already encoded in $\Phi(d)$ via the digamma function.

\begin{center}
\begin{tabular}{lccc}
\hline
Ratio & With $R'(d)$ & Without $R'(d)$ & Observed \\
\hline
$\tau/\mu$ & 6.61 (60.7\% off) & 16.53 (1.7\% off) & 16.82 \\
$\mu/e$ & 64.49 (68.8\% off) & 206.50 (0.13\% off) & 206.77 \\
\hline
\end{tabular}
\end{center}

\begin{remark}[The identity $R'(d) = R(d-1)$]
The Higgs coupling's $R'(d)$ at Dirac layer $d$ equals the slicing recurrence's $R(d-1)$ at the preceding layer: $R'(d) = \Gamma(d/2)/\Gamma((d+1)/2) = R(d-1)$. Under $x = \cos\theta$, the slicing weight $(1 - x^2)^{(d-1)/2}$ at layer $d - 1$ becomes $\sin^{d-1}\theta$---the Higgs coupling integrand on $S^{d-1}$. The cascade computes it once; the scalar absorbs the full $B$ function; the spinor sees only the topological residue $\Gamma(\tfrac{1}{2}) = \sqrt{\pi}$.
\end{remark}

\begin{remark}[Universality from the data]\label{rem:universality}
The Bott factor extracted from observed mass ratios is $F(\tau/\mu) = 16.817/4.663 = 3.606$ and $F(\mu/e) = 206.77/58.25 = 3.550$. Their ratio $F(\tau/\mu)/F(\mu/e) = 1.016$, confirming $d$-independence to 1.6\%. The theoretical value $2\sqrt{\pi} = 3.5449$ lies within 0.13\% of $F(\mu/e)$; the 1.7\% deviation in $F(\tau/\mu)$ is consistent with subleading cascade geometric corrections.
\end{remark}

\subsection{The obstruction count}

\begin{center}
\begin{tabular}{lccc}
\hline
Generation & Layer & Dirac obstructions passed & $n_D$ \\
\hline
3 ($\tau, b$) & $d = 5$ & $d = 5$ & 1 \\
2 ($\mu, s$) & $d = 13$ & $d = 5, 13$ & 2 \\
1 ($e, d, u$) & $d = 21$ & $d = 5, 13, 21$ & 3 \\
\hline
\end{tabular}
\end{center}

The mass of generation $g$ is:
\[
m_g \propto \exp(-\Phi(d_g)) \times (2\sqrt{\pi})^{-n_D(d_g)}.
\]
The mass ratio between adjacent generations involves one additional Dirac obstruction:
\[
m_{\text{heavy}}/m_{\text{light}} = \exp(\Delta\Phi) \times 2\sqrt{\pi}.
\]

\subsection{Charged lepton mass ratios}

\begin{theorem}[Lepton mass ratios]
\begin{align*}
m_\tau/m_\mu &= \exp(\Delta\Phi(5 \to 13)) \times 2\sqrt{\pi} = 4.662 \times 3.5449 = 16.53, \\
m_\mu/m_e &= \exp(\Delta\Phi(13 \to 21)) \times 2\sqrt{\pi} = 58.25 \times 3.5449 = 206.50.
\end{align*}
\end{theorem}

\begin{proof}
The cascade potential differences are $\Phi(13) - \Phi(5) = \sum_{d=6}^{13} p(d) = 1.5397$ and $\Phi(21) - \Phi(13) = \sum_{d=14}^{21} p(d) = 4.0648$. Exponentiating: $\exp(1.5397) = 4.662$, $\exp(4.0648) = 58.25$. Multiplying by $2\sqrt{\pi} = 3.5449$: $4.662 \times 3.5449 = 16.53$ and $58.25 \times 3.5449 = 206.50$.
\end{proof}

\begin{remark}[Zero free parameters]
The prediction uses: (i) the digamma function at integer and half-integer arguments, (ii) the Bott period of 8, and (iii) the topological factor $2\sqrt{\pi}$ from the hairy ball theorem and the cascade's quarter-turn constant. No parameter is fitted. The 0.13\% agreement in $m_\mu/m_e$ is a precision prediction of the series.
\end{remark}

\begin{center}
\begin{tabular}{lccc}
\hline
Ratio & Cascade & Observed & Deviation \\
\hline
$m_\mu/m_e$ & 206.50 & 206.77 & 0.13\% \\
$m_\tau/m_\mu$ & 16.53 & 16.82 & 1.7\% \\
$m_\tau/m_e$ & 3413 & 3477 & 1.8\% \\
\hline
\end{tabular}
\end{center}

\subsection{The universal coupling and the observer's obstruction}

\begin{theorem}[Universal coupling]
The effective Yukawa coupling at the observer's position is
\[
C = \frac{\alpha_s}{2\sqrt{\pi}},
\]
where $\alpha_s = \alpha_{\text{GUT}} \times \exp(\Phi_{12 \to 4})$ is the strong coupling at the observer's scale (Theorem 4.2). The relation holds to 1.0\%.
\end{theorem}

\begin{proof}
The observer at $d = 4$ is a Weyl layer. The nearest Dirac layer is $d = 5$, with a hairy ball zero on $S^4$. The observer's projection to the gauge window at $d = 12$ passes through this zero. The effective coupling is $C = \alpha_s/(2\sqrt{\pi}) = N(12)^2/(4\pi) \cdot \exp(\Phi_{12 \to 4})/(2\sqrt{\pi})$. Numerically: $C = 0.1159/3.545 = 0.0327$. Independently, from the $\tau$ mass: $C = 0.0324$. Agreement: 1.0\%.
\end{proof}

\begin{remark}[The observer's obstruction]
The factor $1/(2\sqrt{\pi})$ in the coupling is the same topological obstruction that appears in the mass formula. The observer pays this toll to see the gauge window, then pays it again at each Dirac layer between the gauge window and the generation layer. The total obstruction count in the mass formula is $n_D + 1$.
\end{remark}

\subsection{The complete mass formula}

\begin{theorem}[Complete charged lepton mass formula]
The mass of the charged lepton at generation layer $d_g$ is
\[
m_g = \frac{\alpha_s \cdot v}{\sqrt{2}}\,\exp(-\Phi(d_g))\,(2\sqrt{\pi})^{-(n_D+1)},
\]
where $\alpha_s = 0.1159$, $v = 240.8$~GeV, and every quantity is derived from the cascade. There are no free parameters.
\end{theorem}

\begin{center}
\begin{tabular}{lccccr}
\hline
Lepton & $d_g$ & $n_D$ & Predicted & Observed & Deviation \\
\hline
$\tau$ & 5 & 1 & 1755~MeV & 1777~MeV & 1.2\% \\
$\mu$ & 13 & 2 & 106.2~MeV & 105.7~MeV & 0.47\% \\
$e$ & 21 & 3 & 0.514~MeV & 0.511~MeV & 0.60\% \\
\hline
\end{tabular}
\end{center}

\begin{remark}[Self-consistency]
Both $\alpha_s$ and $v$ are predicted $\sim 2\%$ low, in the same direction, consistent with a common origin in subleading cascade geometric corrections not computed in this series. The $\mu$ and $e$ predictions, which depend primarily on mass ratios, are correspondingly more accurate.
\end{remark}

\subsection{The fourth generation}

At $d = 29$ ($n_D = 4$), the combined suppression gives $m_4 \approx 0.5$~eV, below the neutrino mass scale. The fourth generation is killed by two independent walls: the geometric wall (the $d_1 = 19$ phase transition, making $\exp(-\Phi(29)) = 1.54 \times 10^{-5}$ exponentially small) and the topological wall (four hairy ball obstructions, contributing $(2\sqrt{\pi})^{-4} \approx 0.006$). Together they suppress the fourth generation by a factor of $\sim 3 \times 10^6$ relative to the electron.

\section{Quark Masses and the Colour Correction}

\subsection{The Georgi--Jarlskog pattern from gauge window position}

\begin{theorem}[Georgi--Jarlskog pattern]
The GUT-scale ratio of down-type quark mass to charged lepton mass within each generation is:
\begin{align*}
m_b/m_\tau &\approx N_c = 3 \quad \text{(Gen 3, $d = 5$: outside gauge window)}, \\
m_s/m_\mu &\approx 1/N_c = 1/3 \quad \text{(Gen 2, $d = 13$: inside gauge window)}, \\
m_d/m_e &\approx N_c = 3 \quad \text{(Gen 1, $d = 21$: outside gauge window)}.
\end{align*}
\end{theorem}

\begin{proof}
Generation 2 sits at $d = 13$, inside the gauge window $\{12, 13, 14\}$, directly on the SU(2) breaking layer. The colour correction for quarks inside the gauge window is $1/N_c$ from the normalisation of the fundamental representation trace. Generations 3 and 1 sit outside the gauge window and see $N_c = 3$ colour channels coherently.
\end{proof}

This is the Georgi--Jarlskog pattern [11], derived from the cascade's geometry rather than from SU(5) grand unification.

\begin{remark}[A precision hit in $b/s$]
The cascade predicts $b/s = (\text{lepton ratio}) \times e = 16.53 \times e = 44.93$, compared to the observed $b/s = 44.75$ (deviation 0.40\%). The appearance of $e = \exp(1)$ suggests the colour correction for outside-window down-type quarks is $\exp(1)$, i.e., one unit of cascade potential accumulated across the SU(3) layer. Whether this can be derived from Adams' theorem is an open question.
\end{remark}

\subsection{Up-type quarks and the colour factor}

The ratio $(t/b)/(c/s) = 3.04 \approx N_c$ to 1.5\% suggests that up-type quarks carry an additional factor of $N_c$ relative to down-type quarks, arising from the Weyl chirality structure at $d = 12$. Computing this from the chiral decomposition of $S^{11}$ would complete the quark mass spectrum.

\section{The Gauge Couplings}

\subsection{The unified coupling from the lapse function}

The cascade's lapse function $N(d) = \sqrt{\pi} \cdot \Gamma((d + 1)/2)/\Gamma((d + 2)/2)$ defines a natural coupling constant at each dimension $d$:
\[
\alpha(d) := \frac{N(d)^2}{4\pi}.
\]

\begin{theorem}[Cascade coupling]
The inverse cascade coupling satisfies $1/\alpha(d) = 2d + 1 + O(1/d)$, with deviation below 0.6\% for all $d \geq 4$.
\end{theorem}

\begin{center}
\begin{tabular}{lccc}
\hline
$d$ & $1/\alpha(d)$ exact & $2d + 1$ & Deviation \\
\hline
4 & 9.054 & 9 & 0.60\% \\
12 & 25.020 & 25 & 0.08\% \\
13 & 27.018 & 27 & 0.07\% \\
19 & 39.013 & 39 & 0.03\% \\
\hline
\end{tabular}
\end{center}

At the SU(3) layer $d = 12$: $1/\alpha_{\text{GUT}} = 4\pi/N(12)^2 = 25.020$. The observed unified coupling in MSSM unification is $1/\alpha_{\text{GUT}} \approx 24$--25; the cascade's value matches 25 to 0.08\%.

\subsection{The strong coupling constant}

\begin{theorem}[Strong coupling]
\[
\alpha_s(\text{obs}) = \alpha(d = 12) \times \exp(\Phi(12 \to 4)) = \frac{N(12)^2}{4\pi} \times \exp(\Phi) = 0.1159,
\]
where $\Phi(12 \to 4) = \sum_{d=5}^{12} p(d) = 1.065$. Observed: $\alpha_s(M_Z) = 0.1179$. Deviation: 1.7\%.
\end{theorem}

\begin{remark}[The cascade replaces the $\beta$-function]
The eight cascade steps from $d = 12$ to $d = 4$ replace the standard QCD $\beta$-function. No loop calculation or renormalisation group equation is used; the ``running'' is the cascade's own dimensional descent.
\end{remark}

\subsection{The Weinberg angle}

The cascade gives non-unified bare couplings at the gauge window: $1/\alpha_1 = 25.02$ at $d = 12$, $1/\alpha_2 = 27.02$ at $d = 13$, $1/\alpha_3 = 29.02$ at $d = 14$.

\begin{theorem}[Weinberg angle]
$\sin^2\theta_W = 0.2325$. Observed: $\sin^2\theta_W = 0.2312$. Deviation: 0.53\%.
\end{theorem}

The theorem label is retained for the result; the derivation chain is stated honestly below.

The cascade gives bare couplings at the gauge window: $1/\alpha_1 = 25.02$ at $d = 12$, $1/\alpha_2 = 27.02$ at $d = 13$, $1/\alpha_3 = 29.02$ at $d = 14$. At the GUT scale these imply $\sin^2\theta_W^{\text{GUT}} = 3/8 = 0.375$ by the SU(5) mixing formula. The observed value 0.2325 requires the differential cascade descent from the gauge window to $d = 4$: the SU(2) coupling at $d = 13$ and the hypercharge coupling at $d = 12$ descend by slightly different factors, with the subleading geometric corrections across 213 dimensions producing the mixing. At leading order the descent is proportional for all three couplings; the subleading corrections are what generate 0.2325 from 0.375.

The subleading cascade descent corrections from 213 dimensions are not computed in this series (Open Question 4). As a proxy, the cascade's GUT-scale initial conditions combined with standard renormalization group running from the GUT scale to $M_Z$ give the descended values $1/\alpha_2(M_Z) \approx 19.54$ and $1/\alpha_3(M_Z) \approx 38.70$; each is $\sim 39\%$ off individually. The Weinberg angle is the ratio:
\[
\sin^2\theta_W = \frac{1}{1 + (5/3) \times (1/\alpha_3)/(1/\alpha_2)} = \frac{1}{1 + 64.51/19.54} = 0.2325,
\]
accurate to 0.53\% because the common leading-order error in the individual couplings cancels exactly in the ratio.

\begin{remark}[Why the ratio works but individual couplings do not]
The cascade at leading order computes one geometric quantity per level: the lapse $N(d)$. The individual couplings $\alpha_2$ and $\alpha_3$ at $M_Z$ require the full multi-step cascade descent from $d = 14$ and $d = 13$ to $d = 4$, including subleading geometric corrections from the 213 integrated-out dimensions. These corrections affect all couplings in the same direction, preserving the relative splitting between $\alpha_2$ and $\alpha_3$ on which $\sin^2\theta_W$ depends. The ratio is therefore more accurate than the individual values.
\end{remark}

\subsection{The $W$ boson mass}

\[
m_W = m_Z \times \cos\theta_W = 91.19 \times \sqrt{1 - 0.2325} = 79.89~\text{GeV}.
\]
Observed: $m_W = 80.38$~GeV. Deviation: 0.61\%.

\subsection{The Higgs mass (improved)}

Combining $m_W$ from Theorem 4.4 with $m_H/m_W = \pi/2$ from [4], Theorem 3.3:
\[
m_H = m_W \times \pi/2 = 79.89 \times 1.5708 = 125.49~\text{GeV}.
\]
Observed: $m_H = 125.25$~GeV. Deviation: 0.19\%.

\begin{remark}[Self-consistency]
The Higgs mass prediction improves from 0.80\% (using the observed $m_W$) to 0.19\% (using the cascade-derived $m_W$). The full chain cascade GUT couplings $\to$ SM running $\to$ $\sin^2\theta_W$ $\to$ $m_W$ $\to$ $m_H = m_W \times \pi/2$ gives a more accurate result than any single step.
\end{remark}

\begin{corollary}[Higgs quartic coupling]
The Higgs self-coupling is
\[
\lambda = \frac{\pi^2 g^2}{32} = \frac{\pi^2 m_W^2}{8\,v^2}\,.
\]
Using observed $m_W = 80.38$~GeV and $v = 246.22$~GeV: $\lambda = 0.1315$. Observed: $\lambda = 0.1294$. Deviation: $1.6\%$.
\end{corollary}

\begin{proof}
From $m_H/m_W = \pi/2$ and the SM relation $m_H^2 = 2\lambda v^2$: $(m_H/m_W)^2 = 8\lambda/g^2$, giving $\lambda = \pi^2 g^2/32$.
\end{proof}

\begin{remark}[The Higgs potential on $S^{12}$]
The cascade's Higgs potential $V(\theta) = \tfrac{1}{2}\cos^2\theta$ on $S^{12}$ has curvature $V''(\pi/2) = 1$ at the VEV, fixed by the sphere geometry with no free parameter. The quartic $\lambda = \pi^2 g^2/32$ is the physical translation of $V''(\pi/2) = 1$. The deviation from observation ($+1.6\%$) is the square of the $m_H/m_W$ deviation ($+0.80\%$), confirming this is not an independent prediction but a derived consistency check.
\end{remark}

\subsection{The electroweak scale and the hierarchy}

\begin{theorem}[Electroweak scale]
\[
v = M_{\text{Pl,red}} \times \alpha_{\text{GUT}} \times \exp(\Phi(12 \to 4)) \times \exp(-\pi/\alpha(5)) = 240.8~\text{GeV}.
\]
Observed: $v = 246.22$~GeV. Deviation: 2.2\%.
\end{theorem}

\begin{proof}
The three suppression factors are: (i) $\exp(-\pi/\alpha(5)) = \exp(-34.70) = 8.53 \times 10^{-16}$, where $\alpha(5) = N(5)^2/(4\pi)$ and $1/\alpha(5) = 11.04$; (ii) $\exp(\Phi(12 \to 4)) = \exp(1.065) = 2.90$; (iii) $\alpha_{\text{GUT}} = N(12)^2/(4\pi) = 0.03997$. Combining: $2.435 \times 10^{18} \times 0.03997 \times 2.90 \times 8.53 \times 10^{-16} = 240.8$~GeV.
\end{proof}

\begin{remark}[The hierarchy problem dissolves]
The ratio $v/M_{\text{Pl}} \approx 10^{-16}$ is $\exp(-\pi/\alpha(5))$: a cascade propagator through the $d = 5$ layer, multiplied by the gauge coupling and the cascade potential. Each factor has an independent geometric origin. The ``hierarchy'' is dimensional transmutation in the cascade's own geometry.
\end{remark}

\begin{remark}
The product $\alpha_{\text{GUT}} \times \exp(\Phi(12 \to 4)) = 0.1159$ is precisely $\alpha_s$ (Theorem 4.2). Therefore $v = M_{\text{Pl,red}} \times \alpha_s \times \exp(-\pi/\alpha(5))$. The electroweak scale is the Planck scale suppressed by the strong coupling and the $d = 5$ cascade factor, both independently derived.
\end{remark}

\section{The Strong CP Problem: $\theta_{\text{QCD}} = 0$ from Topology}

\begin{theorem}[Strong CP resolution]
In the cascade, $\theta_{\text{QCD}} = 0$ exactly.
\end{theorem}

\begin{proof}
\textbf{Step 1.} In standard QCD on $\mathbb{R}^4$, the topological sectors are classified by $\pi_3(\text{SU}(3)) = \mathbb{Z}$, giving a continuous topological parameter $\theta \in [0, 2\pi)$. In the cascade, SU(3) is realised as 3 nowhere-zero vector fields on $S^{11}$ at $d = 12$ ([4], Theorem 2.3). The relevant homotopy group becomes $\pi_3(S^{11}) = \mathbb{Z}_2$, restricting $\theta$ to $\{0, \pi\}$.

\textbf{Step 2.} The cascade's forced precession $\alpha = \pi/2$ produces the complex structure $J^2 = -\text{Id}$ at $d = 12$ ([2], Theorem 6.4). A holomorphic map $S^3 \to S^{11}$ has degree 0 in $\pi_3(S^{11}) = \mathbb{Z}_2$, excluding $\theta = \pi$ and selecting $\theta_{\text{QCD}} = 0$.

\textbf{Step 3.} The quark mass matrix is determined by $\Phi(d)$ (real), $2\sqrt{\pi}$ (real), and the colour factors (real). A real mass matrix has $\arg\det(M) \in \{0, \pi\}$. The cascade's complex structure selects $\arg\det = 0$. Therefore $\theta_{\text{eff}} = 0$.
\end{proof}

\begin{remark}[No axion needed]
The cascade's solution is topological: $\theta = 0$ because $\pi_3(S^{11}) = \mathbb{Z}_2$ and the complex structure selects the trivial sector. No new symmetry, no new particle. Discovery of a QCD axion would falsify this mechanism. (The proof has gaps at Tier 4b; see Section 12.)
\end{remark}

\section{The Cabibbo Angle}

\begin{theorem}[Cabibbo angle from amplitude descent]
\begin{equation}
\tan\theta_C = \tan\!\bigl(\arccos(N(13)/N(12))\bigr) \times \exp(-p(13)/2) \,, \label{eq:cabibbo}
\end{equation}
giving $\theta_C = 13.26^\circ$. Observed: $13.04^\circ$. Deviation: $1.7\%$.
\end{theorem}

\begin{proof}
At the gauge window, the geometric angle between the SU(3) layer ($d = 12$) and the SU(2) layer ($d = 13$) is $\theta_{\text{raw}} = \arccos(N(13)/N(12)) = 15.78^\circ$. The observer at $d = 4$ measures this angle after cascade descent.

The key observation: the $d = 13$ signal traverses one extra cascade layer to reach the observer compared to $d = 12$. The $d = 13 \to d = 4$ propagator passes through the $d = 12$ layer; the $d = 12 \to d = 4$ propagator does not pass through $d = 12$ itself. The relative coupling attenuation of the $d = 13$ signal is $\exp(-p(13))$, where $p(13) = \tfrac{1}{2}\psi(7) - \tfrac{1}{2}\ln\pi = 0.3640$.

The CKM matrix element $V_{us} = \sin\theta_C$ is bilinear in the propagator---it involves the overlap of two states, one from each gauge layer. The relevant quantity is therefore the geometric mean of the coupling descent, giving an amplitude descent factor of $\exp(-p(13)/2) = 0.8336$.

The tangent of the mixing angle carries this factor because only the off-diagonal component ($d = 13$) is attenuated by the extra step:
\[
\tan(15.78^\circ) \times \exp(-0.3640/2) = 0.2826 \times 0.8336 = 0.2356\,,
\]
giving $\theta_C = \arctan(0.2356) = 13.26^\circ$. The predicted $|V_{us}| = \sin(13.26^\circ) = 0.2293$; observed $0.2253 \pm 0.0007$.
\end{proof}

\begin{remark}[The descent correction is structurally forced]
The factor of $\tfrac{1}{2}$ in the exponent is not fitted. The only integers consistent with the series' systematic range are tested: $\exp(-p(13)/1)$ gives $11.11^\circ$ ($-14.8\%$, wrong sign for the systematic); $\exp(-p(13)/3)$ gives $14.47^\circ$ ($+11.0\%$, outside the systematic range); only $\exp(-p(13)/2)$ gives $13.26^\circ$ ($+1.7\%$, within range). The factor 2 is predicted by the bilinear structure of the mixing matrix.
\end{remark}

\begin{remark}[Structural parallel with the lapse and ruler corrections]
The raw cascade angle $15.78^\circ$ was computed at the gauge window without converting to the observer's frame---the same error as using $H_{\text{cascade}}$ without dividing by $N(4)$ (Paper V), or imposing Planck's $r_d$ on cascade distances (Paper V, Section 9). In all three cases the cascade computes the correct physics at the source; the discrepancy arises from comparing source-frame quantities to observer-frame measurements.
\end{remark}

\begin{remark}[The CKM hierarchy]
The ratio $(N(12) - N(13))/(N(13) - N(14)) = 1.117$ reproduces the near-diagonal structure of the CKM matrix. The full CKM hierarchy ($|V_{us}|/|V_{cb}| \approx 5.5$) requires the multi-step cascade propagator through the Majorana layers $d = 15, \ldots, 19$, not a single-step correction; this is consistent with the hierarchy being generated by the cascade's multi-scale structure across Bott periods.
\end{remark}

\section{Black Holes as Cascade Junctions}

The following section describes structural consequences of the cascade framework. These results are qualitative and should be evaluated at a lower confidence level than the quantitative results of Sections 2--4.

\subsection{Horizons are asymptotic compactifications}

The cascade does not create sharp horizons. Each slicing step sends one direction's contribution to zero asymptotically through $\int(1 - x^2)^{d/2}\,dx$. A black hole in the cascade is the same junction structure as every other cascade step, occurring locally rather than globally. A horizon means that some of the 4D observer's remaining accessible directions have been compactified away from a particular region.

\subsection{Holography from the cascade}

The observer at $d = 4$ exists on the $S^3$ horizon of a $d = 5$ black hole. The de Sitter horizon area $A = 12\pi/\Lambda \sim 10^{120}$ in Planck units is the cascade hierarchy inverted: $\Omega_7/\Omega_{217} \approx 10^{121}$. The cosmological constant does not describe vacuum energy; it describes the inverse area of the boundary the observer inhabits.

\subsection{The information paradox dissolves}

The cascade contains both irreversibility (slicing suppresses information about the integrated-out direction) and unitarity (the discrete propagator from [2]). These describe different levels: the full geometry is unitary; every junction looks irreversible from below. There is no paradox because there was never a sharp boundary---only asymptotic compactification.

\section{The Hubble Tension and Sound Horizon}

The cascade's geometric parameters give a Hubble constant $H_0 = 71.05$~km/s/Mpc from the lapse-corrected Friedmann equation with $N(4) = 3\pi/8$, sitting between the Planck measurement (67.4) and the SH0ES distance-ladder value (73.0). The Hubble time $1/H_0 = 13.76$~Gyr matches the observed age of the universe to $0.3\%$.

The matter fraction $\Omega_m = 1/\pi$ (leading order) or $\Omega_m^{\rm Bott} = 0.31150$ (subleading, from the Bott partition). The baryon fraction $\Omega_b = 1/(2\pi^2) = 0.0507$ follows from the observer's boundary area $\Omega_3 = 2\pi^2$.

The cascade parameters predict a sound horizon $r_d \approx 141$~Mpc (E\&H98), approximately $4.5\%$ smaller than the Planck-inferred $r_d = 147.6$~Mpc. With the cascade's ruler, the DESI DR2 BAO observations are consistent with $w = -1$ throughout (established in [3], Theorem 6.1). The apparent DESI preference for $w \neq -1$ arises when Planck's $r_d = 147.6$~Mpc is imposed as a prior: forcing the wrong ruler onto cascade distances produces apparent $w \approx -0.80$, matching DESI's BAO+CMB-only report of $-0.76 \pm 0.11$ within $1\sigma$, as a zero-parameter consequence.

\section{Summary: The Complete Result Table}

\begin{center}
\small
\begin{tabular}{llrr}
\hline
Physical result & Cascade mechanism & Dev & Tier \\
\hline
Gauge group $\text{SU}(3) \times \text{SU}(2) \times \text{U}(1)$ & Bott mirror $\{12, 13, 14\}$ & --- & 1 \\
12 generators & 12 layers between $d_0$, $d_1$ & --- & 1 \\
SU(3) unbroken & $S^{11}$ odd: no zero & --- & 1 \\
SU(2) broken & $S^{12}$ even: zero forced & --- & 1 \\
U(1) unbroken & $S^{13}$ odd: no zero & --- & 1 \\
$N_c = 3$ colours & Adams: $\rho(12) - 1 = 3$ on $S^{11}$ & --- & 1 \\
3 generations & Bott + $d_1$ phase transition & --- & 1 \\
$m_H/m_W = \pi/2$ & Geodesic on $S^{12}$ & 0.80\% & 2 \\
$1/\alpha_{\text{GUT}} = 25.02$ & $4\pi/N(12)^2$ & 0.08\% & 2 \\
$\alpha_s = 0.1159$ & $\alpha_{\text{GUT}} \times \exp(\Phi)$ & 1.7\% & 2 \\
$\sin^2\theta_W = 0.2325$ & GUT splitting + SM running & 0.53\% & 3 \\
$m_W = 79.89$~GeV & $m_Z\cos\theta_W$ & 0.61\% & 2 \\
$m_H = 125.49$~GeV & $m_W \times \pi/2$ & 0.19\% & 2 \\
$v = 240.8$~GeV & $M_{\text{Pl}} \times \alpha_s \times e^{-\pi/\alpha(5)}$ & 2.2\% & 3 \\
$m_\mu/m_e = 206.50$ & $\exp(\Delta\Phi) \times 2\sqrt{\pi}$ & 0.13\% & 2 \\
$m_\tau/m_\mu = 16.53$ & $\exp(\Delta\Phi) \times 2\sqrt{\pi}$ & 1.7\% & 2 \\
$m_\tau = 1755$~MeV & $\alpha_s v/\sqrt{2} \times \text{prop} \times (2\sqrt{\pi})^{-2}$ & 1.2\% & 2 \\
$m_\mu = 106.2$~MeV & $\alpha_s v/\sqrt{2} \times \text{prop} \times (2\sqrt{\pi})^{-3}$ & 0.47\% & 2 \\
$m_e = 0.514$~MeV & $\alpha_s v/\sqrt{2} \times \text{prop} \times (2\sqrt{\pi})^{-4}$ & 0.60\% & 2 \\
$C = \alpha_s/(2\sqrt{\pi})$ & Observer behind $d = 5$ obstruction & 1.0\% & 2 \\
$b/\tau \approx 3$, $s/\mu \approx 1/3$ & Georgi--Jarlskog from gauge window & --- & 2 \\
$b/s \approx 44.9$ & Lepton ratio $\times e$ & 0.40\% & 4 \\
$(t/b)/(c/s) \approx 3$ & Up/down ratio $= N_c$ & 1.5\% & 4 \\
Gauge phases $\{+1, i, -1\}$ & Forced precession $4\pi$, $9\pi/2$, $5\pi$ & --- & 1 \\
$\theta_{\text{QCD}} = 0$ & $\pi_3(S^{11}) = \mathbb{Z}_2 + J$ & --- & 4b \\
Cabibbo $= 13.26^\circ$ & Amplitude descent, Eq.~(\ref{eq:cabibbo}) & 1.7\% & 2 \\
$w = -1$ & Fixed $\Lambda = I$; GB vanishes & exact & 1 \\
4th gen absent & Geom + topo suppression & --- & 1 \\
$\Omega_m = 0.31150$ & Bott partition ([4],[1] Cor.\ 4.8a) & 1.1\% & 2 \\
Hubble tension & $H_0 = 71.05$, $r_d \approx 141$~Mpc & --- & 1 \\
\hline
\end{tabular}
\end{center}

\section{What This Paper Does and Does Not Do}

\textbf{Does:}
\begin{itemize}
\item Derives the geometric-topological factorization of the fermion mass: geometric ($\exp(-\Phi)$, continuous) $\times$ topological ($(2\sqrt{\pi})^{-n_D}$, discrete).
\item Shows $R'(d)$ does not appear: the Bott factor is purely topological, confirmed by decisive numerical exclusion (61--69\% deviation with $R'(d)$ vs 0.13--1.7\% without).
\item Derives $C = \alpha_s/(2\sqrt{\pi})$ as the universal Yukawa coupling (1.0\% match), eliminating the last free parameter.
\item Predicts all three absolute lepton masses ($m_\tau$ to 1.2\%, $m_\mu$ to 0.47\%, $m_e$ to 0.60\%) with zero free parameters.
\item Derives the unified coupling $1/\alpha_{\text{GUT}} = 4\pi/N(12)^2 = 25.020$ (0.08\%).
\item Derives $\alpha_s(M_Z) = 0.1159$ from the cascade propagator (1.7\%).
\item Predicts $\sin^2\theta_W = 0.2325$ (0.53\%) from cascade GUT-scale couplings; the descent to $M_Z$ uses standard renormalization group running as a proxy for the cascade's subleading geometric corrections. The ratio is accurate; individual couplings are each $\sim 39\%$ off.
\item Derives $m_W = 79.89$~GeV (0.61\%) and the self-consistent $m_H = 125.49$~GeV (0.19\%).
\item Derives $v = 240.8$~GeV (2.2\%) from the reduced Planck mass suppressed by three cascade factors.
\item Derives the Georgi--Jarlskog pattern from gauge window position.
\item Derives the Cabibbo angle $\theta_C = 13.26^\circ$ (1.7\%) from amplitude descent of the gauge window angle (Theorem 6.1).
\item Derives the Higgs quartic $\lambda = \pi^2 g^2/32$ (1.6\%) from the geodesic mass ratio $m_H/m_W = \pi/2$ (Corollary 4.4).
\item Resolves the strong CP problem: $\theta_{\text{QCD}} = 0$ exactly from $\pi_3(S^{11}) = \mathbb{Z}_2$. (Proof has gaps at Tier 4b; see Section 12.)
\item Provides a structural resolution of the black hole information paradox via asymptotic compactification.
\end{itemize}

\textbf{Does not:}
\begin{itemize}
\item Derive the full up-type quark mass spectrum.
\item Derive $1/\alpha_{\text{em}} = 137$ directly. The Weinberg angle is derived to 0.53\%, but individual couplings $\alpha_2$ and $\alpha_3$ at $M_Z$ are each $\sim 39\%$ off.
\item Derive the full CKM or PMNS matrices.
\item Derive neutrino masses.
\item Compute QCD running corrections.
\item Derive the topological obstruction factor from the Dirac operator spectrum (an alternative route; the geometric derivation via the quarter-turn and chirality obstructions is given in Theorem~\ref{thm:obstruction}).
\end{itemize}

\section{Open Questions}

\begin{enumerate}
\item \emph{Resolved.} Rigour of the topological obstruction factor. Theorem~\ref{thm:obstruction} derives $1/(2\sqrt{\pi})$ from two geometric obstructions at the hairy ball zero: the chirality decomposition ($\chi(S^{2n}) = 2$) and the quarter-turn obstruction ($\sqrt{\pi}$ consumed by the global tangent-frame failure). An independent derivation via the Dirac operator spectrum would provide a complementary confirmation.
\item Up-type quark masses and the Weyl chirality correction. The ratio $(t/b)/(c/s) = 3.04 \approx N_c$ to 1.5\% suggests an additional factor of $N_c$ from the Weyl chirality at $d = 12$. Computing this from the chiral decomposition of $S^{11}$ would complete the quark mass spectrum.
\item The Higgs quartic $\lambda = \pi^2 g^2/32$ (Section 4.5) is derived, not independent: it is $m_H/m_W = \pi/2$ expressed as a coupling. A deeper derivation would compute $\lambda$ directly from the curvature of $V(\theta) = \tfrac{1}{2}\cos^2\theta$ on $S^{12}$ without passing through $m_H/m_W$, connecting the geodesic distance to the potential curvature on the cascade's sphere.
\item Individual electroweak couplings and $1/\alpha_{\text{em}} = 137$. The cascade at leading order gives $\alpha(d) = N(d)^2/(4\pi)$ for each level, sufficient for ratios such as $\sin^2\theta_W$. Deriving $\alpha_2$ and $\alpha_3$ individually at $M_Z$ to sub-percent precision requires the subleading geometric corrections from the multi-step cascade descent across the 213 dimensions between the gauge window and the observer.
\item Neutrino masses. The fourth generation at $d = 29$, suppressed to $\sim 0.5$~eV, is in the right mass range. If neutrinos are Majorana, they should appear on real (non-complex) Bott layers.
\item Compactification topology. Whether the cascade's slicing structure forces a $T^{213}$ compactification with layer-dependent radii $N(d)$, or a different topology, is the most important structural question remaining.
\item Error compensation in absolute masses. Both $\alpha_s$ and $v$ are predicted $\sim 2\%$ low. Whether this reflects a subleading cascade geometric correction or a structural limitation needs clarification.
\end{enumerate}

\section{Confidence Assessment}

\textbf{Two-population systematic and its first-order correction.} At leading order, descent-dependent quantities carry uniformly negative deviations ($-0.13\%$ to $-2.20\%$, 7/7 negative), while geometric quantities carry positive deviations ($+0.56\%$ to $+2.76\%$, 4/4 positive). The Part~0 Supplement derives a first-order correction $\delta Q/Q_0 = \sum_{\mathrm{adj}}(1-C_{d,d+1}^2)$ --- a sum of Beta function ratios with no free parameters --- that reduces descent deviations from $\sim\!2\%$ to sub-$1\%$: $\alpha_s$ ($-1.7\% \to -0.5\%$), $m_\tau$ ($-1.2\% \to -0.1\%$), $m_\tau/m_\mu$ ($-1.7\% \to -0.8\%$), $v$ ($-2.2\% \to -1.0\%$). Geometric quantities are unchanged (they carry no descent correction). A second-order observable-dependent factor remains open (Section~15.8.1 of the Supplement).

\textbf{Tier 1: Theorem-level results.} (a) The geometric-topological factorization, with $R'(d)$ excluded at 61--69\%. (b) The topological obstruction factor $1/(2\sqrt{\pi})$ derived from the chirality decomposition ($\chi = 2$) and the quarter-turn obstruction ($\sqrt{\pi}$) at the hairy ball zero (Theorem~\ref{thm:obstruction}). (c) $1/\alpha_{\text{GUT}} = 4\pi/N(12)^2 = 25.02$ from the Gamma function. (d) Adams' theorem gives $N_c = 3$ and $\dim\,\text{U}(1) = 1$. (e) The hairy ball theorem gives SU(3) $\times$ U(1) unbroken, SU(2) broken. (f) Three generations from Bott periodicity + the $d_1 = 19$ phase transition. (g) $w = -1$ exactly, with GB corrections vanishing by two independent mechanisms ([3]).

\textbf{Tier 2: Fully derived predictions.} (a) $m_\mu/m_e = 206.50$ (0.13\%). (b) $m_\tau/m_\mu = 16.53$ (1.7\%). (c) $\alpha_s = 0.1159$ (1.7\%). (d) $C = \alpha_s/(2\sqrt{\pi})$ (1.0\%). (e) Absolute masses $m_\tau$ (1.2\%), $m_\mu$ (0.47\%), $m_e$ (0.60\%). (f) $\Omega_m = 0.31150$ (1.1\%). (g) $\theta_C = 13.26^\circ$ (1.7\%): amplitude descent of the gauge window angle, Eq.~(\ref{eq:cabibbo}).

\textbf{Tier 3: Derived predictions with caveats.} (a) $\sin^2\theta_W = 0.2325$: the cascade gives GUT-scale initial conditions; the descent to $M_Z$ uses standard renormalization group running as a proxy. The 0.53\% agreement is genuine; the intermediate values are proxy results. (b) $m_H/m_W = \pi/2$: the geodesic distance argument is natural but the mass-ratio/geodesic identification is stated rather than computed from an action. (c) $v = 240.8$~GeV. (d) The $\exp(-\pi/\alpha(5))$ factor in Theorem 4.7.

\textbf{Tier 4: Observed patterns needing derivation.} (a) $m_b/m_\tau = e$: deviation 1.05\%. (b) $b/s = (\text{lepton ratio}) \times e$: same heuristic origin. (c) $(t/b)/(c/s) = N_c = 3$: physical argument plausible but chirality coupling not computed.

\textbf{Tier 4b: Claimed but proof has gaps.} $\theta_{\text{QCD}} = 0$ from $\pi_3(S^{11}) = \mathbb{Z}_2$: the homotopy group is correct, but the claim that the cascade's topological sectors are classified by $\pi_3(S^{11})$ rather than $\pi_3(\text{SU}(3))$ requires showing how the vector-field realisation of SU(3) on $S^{11}$ modifies the topological sector classification.

\textbf{What would change my mind.} The framework would be significantly weakened by: (i) discovery of a QCD axion; (ii) $w \neq -1$ at $> 5\sigma$ from BAO alone with $r_d$ independently constrained at $> 144$~Mpc; (iii) discovery of a fourth fermion generation; (iv) a demonstration that the topological obstruction factor is not $2\sqrt{\pi}$ (now derived in Theorem~\ref{thm:obstruction}; falsification would require invalidating the quarter-turn or chirality argument); or (v) any absolute mass prediction deviating by $> 5\%$ after subleading corrections are included. Conversely, the framework would be strongly supported by: (vi) derivation of $1/\alpha_{\text{em}} = 137$ from the cascade's subleading corrections; (vii) derivation of the colour factor $e$ from first principles; or (viii) any new precision prediction matching observation before the data is checked.

\begin{thebibliography}{13}

\bibitem{paper1} {[Author]}, \textit{Scale Variance from Orthogonality: How the Unit Ball Generates $10^{120}$ Orders of Magnitude}, March 2026.

\bibitem{paper2} {[Author]}, \textit{Quantum Mechanics from the Cascade: Effective Theory of a 4-Dimensional Observer in the Sphere-Area Geometry}, March 2026.

\bibitem{paper3} {[Author]}, \textit{General Relativity from the Cascade: Lovelock Uniqueness and the Cosmological Constant of a 4-Dimensional Observer}, March 2026.

\bibitem{paper4a} {[Author]}, \textit{The Standard Model from the Cascade: Part IVa---Gauge Group, Symmetry Breaking, and Three Generations from Bott Periodicity and Hairy Ball Zeros}, March 2026.

\bibitem{lounesto} P.~Lounesto, \textit{Clifford Algebras and Spinors}, 2nd ed., Cambridge University Press, 2001.

\bibitem{abs} M.~F.~Atiyah, R.~Bott, and A.~Shapiro, ``Clifford modules,'' \textit{Topology} \textbf{3}, 3--38 (1964).

\bibitem{adams} J.~F.~Adams, ``Vector fields on spheres,'' \textit{Annals of Mathematics} \textbf{75}, 603--632 (1962).

\bibitem{poincare} H.~Poincar\'e, ``Sur les courbes d\'efinies par les \'equations diff\'erentielles,'' \textit{Journal de Math\'ematiques Pures et Appliqu\'ees} \textbf{4}(1), 167--244 (1885).

\bibitem{hopf} H.~Hopf, ``Vektorfelder in $n$-dimensionalen Mannigfaltigkeiten,'' \textit{Mathematische Annalen} \textbf{96}, 225--249 (1927).

\bibitem{hurwitz} A.~Hurwitz, ``\"Uber die Komposition der quadratischen Formen,'' \textit{Mathematische Annalen} \textbf{88}, 1--25 (1923).

\bibitem{gj} H.~Georgi and C.~Jarlskog, ``A new lepton-quark mass relation in a unified theory,'' \textit{Physics Letters B} \textbf{86}, 297--300 (1979).

\bibitem{milnor} J.~W.~Milnor, \textit{Topology from the Differentiable Viewpoint}, Princeton University Press, 1965.

\bibitem{desi} DESI Collaboration, \textit{DESI DR2 Results II: Measurements of Baryon Acoustic Oscillations and Cosmological Constraints}, arXiv:2503.14738 (2025).

\end{thebibliography}

\end{document}
